% Сборка: pdflatex lab04.tex
\documentclass[a4paper,12pt]{article}

\usepackage{cmap}
\usepackage[T1,T2A]{fontenc}
\usepackage[utf8]{inputenc}
\usepackage[english,russian]{babel}

\usepackage{amsmath}
\usepackage{amssymb}
\usepackage[a4paper,margin=2.5cm]{geometry}
\usepackage{enumitem}
% Русские буквы в подпунктах: (а), (б), (в).
\makeatletter
\AddEnumerateCounter{\asbuk}{\@asbuk}{м}
\makeatother
\usepackage{listings}

% Чёрно-белое оформление кода: комментарии курсивом, без цвета.
\lstdefinestyle{bw}{
  language=Python,
  basicstyle=\ttfamily\small,
  commentstyle=\itshape,
  numbers=none,
  frame=single,
  columns=fullflexible,
  keepspaces=true,
  showstringspaces=false,
  breaklines=true,
  upquote=true,
}
\lstset{style=bw}

\newcommand{\pts}[1]{\textbf{(#1)}}

\begin{document}

\begin{center}
  {\Large\bfseries Лабораторная работа 4\par}
  {\large Проверка статистических гипотез\par}
\end{center}

\textbf{Как сдавать работу.}
\begin{enumerate}[itemsep=0.3em]
  \item Сохраните копию ноутбука на свой Диск («Файл $\to$ Сохранить копию на Диске»).
  \item Решите задания в своей копии.
  \item Перед отправкой выберите «Среда выполнения $\to$ Перезапустить и выполнить всё» и убедитесь, что выводы всех ячеек (числа, графики) видны.
  \item В конце ноутбука напишите 2--3 предложения вывода по практической части.
  \item Откройте доступ: «Поделиться $\to$ Все, у кого есть ссылка $\to$ Читатель».
  \item Пришлите преподавателю ссылку на свою копию.
\end{enumerate}

Всего за работу можно получить 10 баллов.

Обозначения: $X_i$~--- элементы случайной выборки (случайные величины), $x_i$~--- их реализация (наблюдаемые значения), $z_q$~--- квантиль уровня $q$ распределения $\mathcal{N}(0,1)$, $\Phi$~--- функция распределения $\mathcal{N}(0,1)$.

\section*{Теоретические вопросы}

\begin{enumerate}
  \item \pts{1 балл} Что такое нулевая гипотеза $H_0$ и альтернативная гипотеза $H_1$, критическая область, ошибки I и II рода, уровень значимости $\alpha$ и мощность критерия?

  Приведите пример из практики разработки (например, проверка, изменилось ли среднее время ответа сервиса после релиза): сформулируйте $H_0$ и $H_1$ и объясните, что в этом примере означают ошибки I и II рода.

  Почему вывод «гипотеза $H_0$ принята» слабее, чем «гипотеза $H_0$ отвергнута»?

  \item \pts{2 балла} Пусть $\mathbf{X}=\left\{X_i\right\}_{i=1}^n$~--- случайная выборка из $\mathcal{N}(\mu,\sigma^2)$, где дисперсия $\sigma^2$ известна. Проверяется $H_0\colon\mu=\mu_0$ против $H_1\colon\mu\ne\mu_0$. Критерий отвергает $H_0$, если $\left|\bar X_n-\mu_0\right|>c$.
  \begin{enumerate}[label=(\asbuk*)]
    \item Какое распределение имеет $\bar X_n$, если $H_0$ верна? Какое распределение имеет статистика $Z=\dfrac{\bar X_n-\mu_0}{\sigma/\sqrt{n}}$?
    \item Найдите порог $c$, при котором вероятность ошибки I рода равна $\alpha$. Запишите правило отклонения $H_0$ через статистику $Z$ (z-критерий).
    \item Пусть $\sigma=1$, $\mu_0=1$, $n=25$, $\alpha=0.05$. При каких значениях выборочного среднего $\bar x$ гипотеза $H_0$ отвергается?
  \end{enumerate}

  \item \pts{1 балл} До релиза среднее время ответа сервиса было $\mu_0=200$ мс, стандартное отклонение известно из мониторинга: $\sigma=40$ мс (будем считать время ответа нормальным). После релиза измерили $n=16$ запросов, среднее время ответа $\bar x=222$ мс.

  Вычислите значение статистики $z$. Изменилось ли среднее время ответа на уровне значимости $\alpha=0.05$? А на уровне $\alpha=0.01$?

  Квантили: $z_{0.975}\approx1.96$, $z_{0.995}\approx2.576$.

\end{enumerate}

\section*{Практические задачи}

\begin{enumerate}
  \item \pts{3 балла} Реализуйте двусторонний z-критерий (вопрос 2): функция по выборке \lstinline|x|, значению $\mu_0$, известному $\sigma$ и уровню значимости $\alpha$ возвращает значение статистики $z$ и решение (\lstinline|True|, если $H_0$ отвергнута). Порог вычисляйте с помощью \lstinline|stats.norm.ppf|.

  Выберите $\mu$ и $\sigma$ (например, время ответа $\mu=200$, $\sigma=40$), сгенерируйте с помощью \lstinline|np.random.default_rng(seed)| выборку объема $n=30$ из $\mathcal{N}(\mu,\sigma^2)$ и проверьте $H_0\colon\mu=\mu_0$ на уровне $\alpha=0.05$ для двух значений: $\mu_0=\mu$ (гипотеза верна) и $\mu_0=\mu+\sigma/2$ (гипотеза неверна).

  Найдите порог критерия моделированием: сгенерируйте $M=100000$ значений $Z\sim\mathcal{N}(0,1)$ и вычислите квантиль уровня $1-\alpha$ величины $|Z|$ (\lstinline|np.quantile|). Сравните с $z_{1-\alpha/2}$ из \lstinline|stats.norm.ppf|.
\begin{lstlisting}
def z_test(x: np.ndarray, mu0: float, sigma: float, alpha: float) -> tuple[float, bool]:
    # your code
    return z, reject
\end{lstlisting}

  \item \pts{3 балла} Проверьте, что вероятность ошибки I рода z-критерия равна $\alpha$. Для $\mu_0=\mu$ (гипотеза $H_0$ верна) сгенерируйте $M=10000$ выборок объема $n=20$, для каждой вычислите статистику $z$ функцией \lstinline|z_test| и найдите долю выборок, в которых $H_0$ отвергнута, для $\alpha=0.01,\ 0.05,\ 0.1$.

  Постройте гистограмму значений $z$ (с \lstinline|density=True|), плотность $\mathcal{N}(0,1)$ и вертикальные линии порогов $\pm z_{0.975}$. Какое распределение, судя по гистограмме, имеет $z$, если $H_0$ верна? Как это объясняет, что доля отвержений равна $\alpha$?

\end{enumerate}

\section*{Вывод}

Напишите 2--3 предложения вывода по практической части: отвергается ли $H_0$ в задаче 1 для верной и неверной гипотезы, какова доля ложных отклонений при верной $H_0$ и как распределена статистика $z$.

\end{document}
